\section{Discussion}
\label{sec:discussion}

The three datasets converge on a single structural claim and a set of qualifications around it. We synthesize them here in the order of the contribution, then recap the framing discipline that keeps the claim honest. No new figures are introduced beyond those already reported in Sections~\ref{sec:results} and elsewhere; this section interprets them.

\subsection{Constitutive enforcement subsumes corrective enforcement}
\label{subsec:constitutive-subsumes}

The central result is not that GCD blocks attacks --- a correctly wired post-parse allowlist (\texttt{D}) blocks the same \emph{effect} on the same policy, because \texttt{D} and \texttt{C} enforce the same language \(L(G_s)\) through the same \texttt{tantalus\_grammar::allowlist\_verdict} (\texttt{D == C} by construction). The result is that \texttt{C} and \texttt{D} differ in \emph{enforcement completeness}: where the policy sits relative to the generative act. \texttt{C} is constitutive --- the policy \emph{is} the output space, so the out-of-scope action is ungenerable and never exists as an artifact on any surface (output buffer, logs, telemetry, stream, KV cache). \texttt{D} is corrective --- an unbounded, default-allow generator produces the forbidden action in full, after which a separate gate inspects and rejects it.

The empirical signature of that difference is \emph{emission}, the per-trial count of times the unauthorized artifact was actually produced. Across both victims the signature is the same. On the weak Qwen3-1.7B workhorse, the behavioral generators behind the allowlist gate (\texttt{a1\_d\_r3}--\texttt{a4\_d\_r3}) each produced \num{104615}--\num{120140} out-of-scope calls (emission), drove \num{111528}--\num{139612} gate rejections, and stranded 44--46 trials at the retry-3 budget; the three GCD arms (\texttt{c\_d\_r3}, \texttt{c\_closed\_d\_r3}, \texttt{c\_guided\_d\_r3}) recorded \emph{0} emission, \emph{0} gate rejections, and \emph{0} \texttt{availability\_failure} --- the gate is dead weight, because the grammar already enforced the policy at the sampler. On the capable Mistral-119B anchor the same shape holds: \texttt{c} emits 0 while the corrective \texttt{D} arms produce \num{2317}--\num{3878} calls and pay an availability cost (\texttt{d\_r0} strands all 3878; retry recovers to 288 at \texttt{d\_r1} and 6 at \texttt{d\_r3}, at \numrange{1.63}{1.72} average attempts, max 5). This is the produce-then-catch tax, and it is precisely the failure mode positive security exists to escape, reintroduced one layer down: the dangerous artifact is produced (a confidentiality and forensics surface), a coverage obligation attaches to the gate (it must fire on every output path, consumer, and refactor or be bypassed), and a time-of-check/time-of-use window opens. GCD has none of these because there is nothing to catch.

We state the boundary of the contribution carefully. \texttt{C} is not \emph{more secure} than a perfectly wired allowlist on the blocked effect --- they enforce the same positive policy. The difference is that GCD's positivity is enforcement-deep (positive security all the way down to the generative act) while the allowlist's is policy-deep only; the same post-parse check is therefore \texttt{C}'s recommended corrective backstop in deployment, not a dominated alternative (Section~\ref{subsec:framing}). The cost ordering \emph{among} the behavioral variants on the 1.7B inherits the sub-floor-victim caveat (Section~\ref{subsec:incompleteness}); the load-bearing, victim-independent result is the qualitative 0-vs-\(>\!0\) gate inertness.

A caveat on the token-cost axis follows directly. The pre-registered H3 decision rule is conjunctive: \texttt{C} emission \(=0\) \emph{and} \texttt{C} attack-trial output-token cost strictly below \texttt{D}'s. On the 1.7B both conjuncts hold (\texttt{C} 119 output tokens vs.\ \texttt{D} 261--289). On Mistral the emission conjunct carries decisively but the token-cost conjunct fails: blind-\texttt{C}'s 208 output tokens are among the leanest yet not strictly below every \texttt{D} variant (\texttt{d\_r1} 203). A capable model writes a substantive free \texttt{respondToUser} message, narrowing the token gap that was large on the weak victim. So the formal, conjunctive H3 rule is met only on the 1.7B; on Mistral the qualitative gate inertness and the validity dividend (below) carry, while the token ordering is victim-dependent. We do not claim H3-as-written is satisfied on both victims.

\subsection{The validity dividend is a model-independent reliability property}
\label{subsec:validity-dividend}

Beyond security, GCD buys reliability that the model does not supply on its own. Across both victims the GCD arms produced a well-formed, executable, in-policy tool call far more often than native free generation did: \texttt{c\_d\_r3} 98.9\%, \texttt{c\_closed\_d\_r3} 99.9\%, \texttt{c\_guided\_d\_r3} 97.8\% valid output on the 1.7B (vs.\ \numrange{89.4}{93.9}\% for the behavioral arms), and \texttt{c} 99.8\% valid on Mistral against only \(\sim\!59\%\) for native free generation. The capable model emits a valid tool call barely more than half the time --- the rest is tool-call-as-text, spin-out, or truncation that never executes --- and the grammar eliminates that failure class by construction. This is a reliability property of the enforcement mechanism, not of the model: GCD guarantees a structurally valid in-policy call the native model fails to produce a substantial fraction of the time, independent of model scale or competence.

Two qualifications keep this honest. First, ``0 invalid outputs'' is \emph{not} literally met on vLLM: under the Kleene-\texttt{*} free-string rules the weak 1.7B occasionally rambles \texttt{respondToUser} or \texttt{searchFiles.query} to \texttt{max\_tokens} and truncates mid-string, leaving an unparseable turn. The residual is largest for \texttt{c\_guided\_d\_r3} ($\approx\!2.2\%$ invalid, above the pre-registered floor; recorded as a deviation), whose terminal reply is still a free respond-only string. It fails safe --- a truncated grammar output yields no executable call, never a wrong action, so 0 win / 0 emission still holds on every such trial --- and remains far better than the \numrange{8}{10}\% parse failures of native generation, but it is not literal zero, and the fix path (a closed canned respond-only terminal, as \texttt{C+} already applies to \texttt{respondToUser}) is named rather than claimed. Second, the validity gain does not make a small model capable: intelligence lives in the corpus and the selection, not in the grammar. GCD makes the model's \emph{output} valid and in-policy; it does not make the model frontier.

\subsection{Negative-security incompleteness, made visible}
\label{subsec:incompleteness}

On the capable Mistral anchor the behavioral stack behaves exactly as Galloway's negative/positive dichotomy predicts \cite{galloway2020}. The in-band controls are inert against injection: the defensive prompt A1 (56.50\% bypass) does not improve on Control (56.77\%) --- the A1\,\(\le\)\,Control manipulation check holds --- and the input classifier A2 (56.85\%) is likewise useless, both because they live in the same input channel the injected skill occupies, so the injection overrides them. The output filter A3/A4 does the heavy lifting (2.15\% / 1.82\% bypass) but \emph{never reaches 0}: A4's CI lower bound of 1.43\% is strictly above zero, satisfying H1.

The staging funnel localizes \emph{why} A3/A4 still leak. The attack reaches the \texttt{fetchURL} sink in \(\sim\!57\text{--}58\%\) of trials in \emph{every} arm, A3/A4 included; A1 and A2 do not move the funnel at all, and A3/A4 leave the dangerous \emph{emission unchanged}, blocking only at the final output stage. A4's 1.82\% is therefore ``the attack fired \(\sim\!58\%\) of the time and the credential denylist caught all but 73,'' not ``the attack rarely fired.'' The 73 (\texttt{a4}) and 86 (\texttt{a3}) survivors are exactly the credential-signature-free exfils --- the TicketAssist business-data class H1 specifically predicted --- which carry no pattern for a denylist to match. This is the produce-then-catch tax surfacing \emph{within} a behavioral arm: a post-generation credential denylist is itself corrective and incomplete by construction, so the unauthorized \texttt{fetchURL} is generated regardless of the defense, and any payload outside the signature set passes. The contrast with \texttt{C}'s constitutive 0 emission is the whole point.

\subsection{The un-abliteratable sampler}
\label{subsec:un-abliteratable}

The abliterated-adversary steelman isolates \emph{where} each control lives. Behavioral safety lives in the weights, and abliteration \cite{arditi2024} surgically removes the refusal directions that A1--A4 implicitly rely on --- defeating the entire behavioral stack by construction. Under a top-authority malicious prompt the abliterated Qwen3.5-27B is demonstrably willing and able, bypassing 97.85\% (1957/2000) in the free cell. Constitutive enforcement, by contrast, lives at the decode boundary, in the sampler, not in the weights: the grammar mask drives every grammar-invalid token's probability to zero \emph{before} sampling, so the same model under the same prompt is driven to 0.00\% (0/2000) bypass with the grammar as the sole independent variable. The transformer's logits --- intent --- remain uncontrolled and may, under injection, place arbitrarily high probability on a forbidden token; the \emph{generated tokens} are nonetheless constrained.

We present this as a rebuttal exhibit, not main statistics, and bound it precisely. \texttt{c\_sealed}'s 0\% is partly true-by-construction --- the forced denylist trajectory seals the exfil sink --- so its empirical content is the paired open cell (the adversary is real at 97.85\%) and the engagement evidence: in 96.8\% (1937/2000) of sealed trials the model \emph{provably ran its skill's forced sensitive-artifact read} before reaching the sealed safe-fetch and the verbatim taunt (AWS-credential file in 542, SSH key 259, cross-user chat 517, ticket business-data 619). That answers the two natural objections --- ``the model just refused'' (it cannot, and bypasses 97.85\% unconstrained) and ``it never read the bad content'' (96.8\% provably did). It is engagement evidence, not a confidentiality result: GCD here \emph{permits} the read and blocks only the exfil action; the AWS-credential claim holds for the 542 ApiKey/Infra trials, not for all 1937. Protecting a secret through free-text channels is explicitly Paper-2 scope.

\subsection{Cost and economics}
\label{subsec:cost}

Token generation is not free, and the constitutive/corrective distinction has a direct economic reading. A corrective gate must let the generator spend the tokens to produce the forbidden artifact, then discard them; a constitutive grammar never spends them. On attack trials this shows up as the GCD arms generating $\approx\!40\text{--}55\%$ of the behavioral output tokens on the 1.7B (119--149 vs.\ 261--289), with the gap narrowing on a capable model whose blind-\texttt{C} reply is substantive (208 vs.\ \texttt{D}'s 203--270). We frame this as a layer on top of the security claim, not a substitute, and we make no single-axis ``cheaper'' claim: \texttt{C} carries a per-token masking cost that is engine-dependent (a CPU mask cliff on \texttt{llama.cpp} \cite{gerganov_llamacpp}; at-par on vLLM with xgrammar \cite{dong2024xgrammar}, which has no jump-forward but also no cliff), so \texttt{C} is not ``the slow arm'' --- the avoided regeneration is a real but bounded saving, reported against an honest per-engine latency picture. The third C-vs-D signature, time-to-valid-action measured by \texttt{attempts}, reinforces the structural reading: \texttt{C} is one-shot (one attempt by construction, no gate round-trip), while \texttt{D} serializes 1.0\,\(\to\)\,1.72 average attempts plus a rejection round-trip per retry.

\subsection{Framing discipline: what is and is not load-bearing}
\label{subsec:framing}

Security at scale is the headline. The contribution ports Galloway's positive/structural-vs-negative/behavioral comparison \cite{galloway2020} into the LLM-agent domain: the same attack corpus, the same agent, only the control changes, scored by \emph{effect} with Clopper-Pearson intervals \cite{clopper1934}, at a scale roughly 26--27\(\times\) ASB-2024's 90{,}000 test cases \cite{zhang2024asb}, with 852{,}000 GCD attack trials at 0 bypass. The soundness proof and the cost axis are layers \emph{on top} of that headline, not stand-ins for it.

Three disciplines keep the empirical claim from collapsing into a tautology. First, \texttt{C}'s 0\% is load-bearing because it is run at scale, scored by effect, and measured \emph{through} the unconstrained free-text surface (\texttt{searchFiles.query}, \texttt{respondToUser.message} remain free): bypass is ``did the marker reach an out-of-scope sink,'' not ``is the output a member of \(L(G_s)\).'' The grammar does not co-define the threat model; the win detector is independent of grammar membership. This is the in-band/out-of-band split made operational --- the behavioral controls fail because the attacker shares their channel, while the grammar is compiled from authenticated identity and applied at the sampler, nowhere the attacker's tokens can reach. Second, the claim is scoped to \emph{action integrity over enumerable-sink tools}; free-text confidentiality through channels that must remain free-text is stated up front as out of scope (Paper-2), consistent with the broader concern that injection through trusted data channels is a structural, not incidental, exposure \cite{greshake2023}. Third, the corpus is selected-on-outcome --- attacks known effective against the undefended agent, the live-malware analogue --- so it is the population of \emph{capable} attacks, not a representative sample of the threat space; bypass rates are conditional on attack potency against the no-defense replay baseline ($\sim\!51.8\%$), and the large N tightens intervals without buying generalization.

Two further boundaries scope what the experiment can and cannot underwrite. The ``eliminates by construction'' claim is carried by the soundness theorem, which quantifies over all logit distributions --- it holds against a fully injection-compromised model because it never assumes the logits are benign, giving
\[
\text{emitted} \in L(G_s) \subseteq \mathrm{Authorized}(s),
\]
and no finite sample could license a class-level elimination claim that this proof does \cite{willard2023}. The experiment instead shows that behavioral controls genuinely leak at scale and that the implementation matches the theory; proof and evidence stay in separate lanes. And the empirical claim itself is a feasibility claim on a single subject system (Tantalus), not a population-level universality claim. The trusted computing base is correspondingly explicit: the guarantee holds iff \(L(G_s) \subseteq \mathrm{Authorized}(s)\) --- i.e.\ iff the compiled allowlist is correct --- which is why we recommend the post-parse check as an independent acceptance oracle and executor re-validation backstop, a Condition-\texttt{D}-class corrective control deployed as defense-in-depth that complements \texttt{C} rather than competing with it.

Finally, external validity remains the principal open boundary, and we do not overstate it. RQ1, the formal H2 contrast, and the blind-\texttt{C} utility point estimate are demonstrated on one capable non-Qwen victim (Mistral-119B), but the headline external-validity claim is not yet closed: the pre-registration ties it to \(\ge 3\) distinct pretraining lineages, and we are at 2 of 3 (Qwen and Mistral). The C=0 invariant additionally reproduced across three inference engines and three grammar backends, which supports engine-independence as a robustness note --- not as three co-equal anchors. The registered powered 2pp non-inferiority TOST for RQ4 was run on neither victim, so we claim at most a point-estimate ordering (blind-\texttt{C} legit-success above Control on the capable model, an 11-percentage-point deficit on the weak one), and we attribute the deficit to weak-model competence rather than to the grammar, with \texttt{c\_guided} (RQ6: 100.0\% legitimate-task success, 0 bypass, 0 deflection on the 1.7B) the universal resolution. These are deferrals, not contradictions, and they bound the claim to exactly what the three datasets support.
